% ECA_21.tex
\documentclass[a4paper,11pt]{ltjsarticle}
\usepackage{amsmath,amssymb,amsthm}
\usepackage{mathtools}
\usepackage{geometry}
\geometry{margin=25mm}

\title{ECA_21 $I$ の選択閉包性と $B(i)$ の役割}
\author{}
\date{}

\newtheorem{definition}{定義}
\newtheorem{proposition}{命題}
\newtheorem{remark}{注意}

\begin{document}
\maketitle

\section{目的}

ECA_20 では、可算無限公理候補族
\[
\mathcal{R}^*=\{r_n\mid n\in\mathbb{N}\}
\]
に対して
\[
\forall A\subseteq\mathcal{R}^*,
\qquad
\mathcal{A}_{\mathrm{ZFC}}\cup A\in I
\]
という局所的な選択閉包条件を検討した。

本稿では、この条件が ECA の既存構造
\[
I\subseteq\mathcal{P}(\Ax),
\qquad
B(i)\subseteq i\subseteq\Ax
\]
から自然に導出される可能性を検査する。

特に、
\[
B(i)
\]
が単なる部分集合として定義されているだけなのか、それとも公理選択集合 $i$ の構成規則に関する条件を与えるものなのかを区別する。

\section{現行の基本関係}

公理候補集合を
\[
\Ax
\]
とし、公理選択空間を
\[
I\subseteq\mathcal{P}(\Ax)
\]
とする。

各
\[
i\in I
\]
について
\[
B(i)\subseteq i\subseteq\Ax
\]
を満たす。

また、Solver
\[
s=(\mathcal{A}_s,\mathcal{R}_s)
\]
について
\[
\mathcal{A}_s\in I
\]
を要求する。

したがって、
\[
B(i)
\]
は少なくとも形式上は、選択された公理集合 $i$ の内部にある部分集合である。

\section{$B(i)$ から全部分集合閉包を導けるか}

現在の関係
\[
B(i)\subseteq i
\]
だけから、
\[
A\subseteq i
\Longrightarrow A\in I
\]
を導くことはできない。

そもそも $B(i)$ の定義は
\[
B(i)\subseteq i
\]
であり、
\[
B(i)=\mathcal{P}(i)
\]
でも
\[
B(i)=\{A\subseteq i\mid A\in I\}
\]
でもない。

したがって、現行の定義だけでは $B(i)$ によって $I$ の閉包性を表現することはできない。

\begin{proposition}
\[
I\subseteq\mathcal{P}(\Ax),
\qquad
B(i)\subseteq i
\]
だけでは、
\[
\forall A\subseteq\mathcal{R}^*,
\quad
\mathcal{A}_{\mathrm{ZFC}}\cup A\in I
\]
は導出できない。
\end{proposition}

\begin{proof}
例えば
\[
I=\{\mathcal{A}_{\mathrm{ZFC}}\}
\]
とし、
\[
B(\mathcal{A}_{\mathrm{ZFC}})
=
\varnothing
\]
とすれば、
\[
B(i)\subseteq i\subseteq\Ax
\]
は成立する。

一方、$\mathcal{R}^*\neq\varnothing$ なら、一般に
\[
\mathcal{A}_{\mathrm{ZFC}}\cup A
\notin I
\]
となる $A\subseteq\mathcal{R}^*$ が存在する。

よって、現在の包含関係だけでは選択閉包性は導けない。
\end{proof}

\section{$B(i)$ を選択可能部分族として読む場合}

別の定義として、
\[
B(i)
=
\{A\subseteq i\mid A\text{ が許容される}\}
\]
のように、$B(i)$ を $i$ 内部の選択可能部分集合族として定義することは考えられる。

この場合には
\[
B(i)\subseteq\mathcal{P}(i)
\]
となる。

さらに
\[
B(i)=\mathcal{P}(i)
\]
を仮定すれば、$i$ の任意の部分集合が $B(i)$ に属する。

しかし、これは現在の
\[
B(i)\subseteq i
\]
とは別の定義である。

したがって、この解釈を ECA に導入する場合は、既存の $B(i)$ の意味を変更することになるため、現時点では導入せず、別案として扱う。

\section{$I$ の局所閉包性を直接定義する場合}

$B(i)$ を変更せずに、$I$ の閉包性を直接記述する方法がある。

対象となる
\[
\mathcal{R}^*\subseteq\Ax
\]
と
\[
\mathcal{A}_{\mathrm{ZFC}}\subseteq\Ax
\]
に対して、
\[
\boxed{
\forall A\subseteq\mathcal{R}^*,
\quad
\mathcal{A}_{\mathrm{ZFC}}\cup A\in I
}
\]
を直接条件として置く。

この方法では
\[
B(i)
\]
の意味を変更する必要がない。

また、
\[
I=\mathcal{P}(\Ax)
\]
という大域的な条件も不要である。

したがって、現行構造をできるだけ保存するという観点では、この記述が最も直接的である。

\section{$B(i)$ と局所閉包条件の関係}

仮に
\[
i_A=\mathcal{A}_{\mathrm{ZFC}}\cup A
\]
と置く。

ECA_20 の局所閉包条件は
\[
\forall A\subseteq\mathcal{R}^*,
\quad i_A\in I
\]
である。

一方、
\[
B(i_A)\subseteq i_A
\]
は各 $i_A$ が選択された後に成立する包含関係である。

したがって、論理的には

\[
\boxed{
i_A\in I
\quad\text{と}\quad
B(i_A)\subseteq i_A
}
\]

は別の情報である。

前者は「その公理集合が選択空間 $I$ に属する」ことを表し、後者は「その公理集合に対応する $B$ が $i_A$ の部分集合である」ことを表す。

このため、現在の定義では $B(i)$ を用いて Stage 2 の選択閉包性を導く必要はない。

\section{有限閉包と全部分集合閉包}

$B(i)$ に何らかの閉包性を追加する場合にも、有限の場合と全部分集合の場合を区別する必要がある。

例えば
\[
A,B\subseteq\mathcal{R}^*
\]
について
\[
i_A,i_B\in I
\]
から
\[
i_{A\cup B}\in I
\]
を導く二項閉包性を考えることができる。

しかし、このような有限的な閉包性だけでは
\[
\forall A\subseteq\mathcal{R}^*,
\quad i_A\in I
\]
を直ちに導けない。

特に可算無限集合
\[
A=\{r_n\mid n\in J\},
\qquad J\subseteq\mathbb{N}
\]
については、有限近似
\[
A_n=A\cap\{r_0,\ldots,r_{n-1}\}
\]
から
\[
i_A\in I
\]
を得るための極限・可算閉包に相当する条件が別途必要になる。

\section{選択閉包を ECA の構造条件とする場合}

ECA の基本構造を変更することなく、選択空間 $I$ に対する追加条件として
\[
\mathsf{Closure}(\mathcal{R}^*,\mathcal{A}_{\mathrm{ZFC}})
\]
を導入し、
\[
\mathsf{Closure}(\mathcal{R}^*,\mathcal{A}_{\mathrm{ZFC}})
\iff
\forall A\subseteq\mathcal{R}^*,
\quad
\mathcal{A}_{\mathrm{ZFC}}\cup A\in I
\]
と定義することができる。

この場合、
\[
\mathsf{Closure}
\]
は $B(i)$ の定義ではなく、$I$ に対する追加の構造条件となる。

この分離には、

\[
\operatorname{Ax}(U)
\quad\longrightarrow\quad
\mathcal{R}^*
\]

と

\[
\mathcal{R}^*
\quad\longrightarrow\quad
I
\]

を別々に検証できるという利点がある。

\section{連続体濃度への影響}

\[
|\mathcal{R}^*|=\aleph_0
\]
および
\[
\mathsf{Closure}(\mathcal{R}^*,\mathcal{A}_{\mathrm{ZFC}})
\]
に加えて
\[
\mathcal{R}^*\cap\mathcal{A}_{\mathrm{ZFC}}=\varnothing
\]
を仮定すると、
\[
\Gamma:
\mathcal{P}(\mathcal{R}^*)\hookrightarrow I,
\qquad
\Gamma(A)=\mathcal{A}_{\mathrm{ZFC}}\cup A
\]
が得られる。

したがって
\[
|I|\ge2^{\aleph_0}.
\]

しかし、これは依然として
\[
|\mathcal Z|\ge2^{\aleph_0}
\]
を意味しない。

Stage 3 において
\[
\mathcal T(s_A)\cong_{\mathrm{Th}}\mathrm{ZFC}
\]
をすべての $A\subseteq\mathcal{R}^*$ について確認する必要がある。

\section{現段階で採用すべき構造}

現時点では、
\[
B(i)\subseteq i
\]
という既存の定義を変更せず、

\[
\boxed{
\forall A\subseteq\mathcal{R}^*,
\quad
\mathcal{A}_{\mathrm{ZFC}}\cup A\in I
}
\]

を独立した Stage 2 の検証条件として扱うのが最も整合的である。

これにより、

\[
\operatorname{Ax}(U)
\to
\mathcal{R}^*
\]

は ECA_19、

\[
\mathcal{R}^*
\to
I
\]

は本稿、

\[
I
\to
\mathcal Z
\]

は Stage 3 として分離できる。

\section{結論}

現在の ECA における
\[
I\subseteq\mathcal{P}(\Ax),
\qquad
B(i)\subseteq i
\]
から、$\mathcal{R}^*$ の全部分集合に対する選択閉包性を導くことはできない。

特に、
\[
B(i)\subseteq i
\]
は、それ自体では $I$ の閉包性を規定していない。

したがって、現段階では
\[
B(i)
\]
の意味を変更して閉包性を担わせるより、

\[
\boxed{
\forall A\subseteq\mathcal{R}^*,
\quad
\mathcal{A}_{\mathrm{ZFC}}\cup A\in I
}
\]
を独立した選択構造上の条件として検査する方が、既存の ECA 構造を保ったまま議論できる。

次の検討では、この条件をさらに
\[
I
\]
そのものの定義、すなわち「何を ECA における許容可能な公理選択とするのか」という問題まで戻して調べる必要がある。

\end{document}
